\subsubsection{Směr střelby a automatické zaměření na cíl}

Střelba se automaticky zaměřuje na bowlingovou kouli ve scéně. To bylo zvoleno proto, aby hráč mohl využít obě ruce na klávesnici pro pohyb a další schopnosti. Výpočet směru projektilu využívá normalizovaný vektor mezi pozicí hráče a pozicí koule.

Vzorec pro výpočet směru je \texttt{(ballPosition - playerPosition).normalized()}. Tento vektor udává směr od hráče k cíli a jeho normalizace zajistí, že projektil má konstantní rychlost bez ohledu na vzdálenost.

Projektil se instancuje na pozici \texttt{firePoint}, což je bod na hráči odkud vylétá projektil. Rotace projektilu se nastaví podle směru letu pomocí \texttt{Mathf.Atan2} -- tento úhel se převede na stupně a nastaví se jako rotace kolem osy Z. Toto zajišťuje, že projektil je vizuálně orientován ve směru pohybu.

\begin{lstlisting}[language=csharp,caption={Listing 12: Metoda Shoot pro vytvoření projektilu}]
void Shoot()
{
    Vector2 direction = (enemyTransform.position - firePoint.position).normalized;

    GameObject bullet = Instantiate(bulletPrefab, firePoint.position, Quaternion.identity);
    Rigidbody2D rb = bullet.GetComponent<Rigidbody2D>();
    rb.linearVelocity = direction * bulletSpeed;

    float rotZ = Mathf.Atan2(direction.y, direction.x) * Mathf.Rad2Deg;
    bullet.transform.rotation = Quaternion.Euler(0, 0, rotZ);
}
\end{lstlisting}
